\section{A fixed number of disjoint cycles at high path density}
\label{sec:fixed-cycle-packings}

The thin-layer encoding extends to any fixed number of cycle components.
Its polynomial factor is the number of ways to split a long sequence of ordinary layers.
This gives an optimal polynomial order for each fixed number of components.

\begin{lemma}[A fixed number of almost disjoint intervals]
\label{lem:fixed-interval-tuple-sum}
Fix integers $j\ge1$ and $B\ge0$.
There is a constant $c_{j,B}$ such that
\[
 \sum_{\substack{I_1,\ldots,I_j\text{ nonempty integer intervals in }[0,D]\\
 |I_a\cap I_b|\le B\ (a\ne b)}}
 2^{\frac12\sum_{a=1}^j|I_a|}
 \le c_{j,B}(D+1)^{j-1}2^{D/2}
\]
for every integer $D\ge0$.
The interval tuples are ordered.
\end{lemma}

\begin{proof}
Put $u=\sqrt2$ and $N=D+1$.
First consider $q$ pairwise disjoint intervals in their left-to-right order.
Let their positive lengths be $\ell_1,\ldots,\ell_q$.
Let $g_0,\ldots,g_q$ be the nonnegative gaps, including the two outer gaps.
Then
\[
 \sum_i\ell_i+\sum_i g_i=N.
\]
For a fixed gap vector, there are at most $N^{q-1}$ possible length vectors.
Thus the weighted sum over these disjoint intervals is at most
\begin{equation}\label{eq:disjoint-interval-tuple-bound}
 N^{q-1}u^N\sum_{g_0,\ldots,g_q\ge0}u^{-\sum_i g_i}
 =\frac{u}{(1-u^{-1})^{q+1}}N^{q-1}u^D.
\end{equation}

Now sort a general tuple by its left endpoints, at a cost of at most $j!$.
For each interval, remove the part covered by earlier intervals.
Because all earlier left endpoints are no larger, the removed part is an initial segment.
The remaining interval is empty or is a single interval.
The nonempty remainders are pairwise disjoint and occur in order.

The $i$-th interval loses at most $(i-1)B$ positions.
Indeed, its overlap with the union of earlier intervals is at most the sum of its pairwise overlaps.
The total removed length is therefore at most $B\binom j2$.
At least the first remainder is nonempty.
Let $q\ge1$ be the number of nonempty remainders.

For a fixed nonempty remainder, its original interval is recovered by restoring an initial segment of length at most $(j-1)B$.
There are at most $(j-1)B+1$ possibilities.
An empty remainder came from an interval of length at most $(j-1)B$.
There are at most $((j-1)B+1)N$ possibilities for such an interval.
Choosing which remainders are empty costs only a factor depending on $j$.
Their possible positions therefore contribute at most a constant times $N^{j-q}$.
The weight lost in trimming contributes at most $u^{B\binom j2}$.

Apply \eqref{eq:disjoint-interval-tuple-bound} to the $q$ nonempty remainders.
The resulting power is $N^{j-q}N^{q-1}=N^{j-1}$.
All other factors depend only on $j,B$.
Summing over $1\le q\le j$ proves the lemma.
\end{proof}

\begin{theorem}[Fixed-component cycle packing bound]
\label{thm:fixed-component-high-path-packings}
Fix $j\ge1$ and $0<\varepsilon\le1/2$.
There is a constant $K_{j,\varepsilon}$ such that every normalized two-pole $N$ of rank $b$ satisfying
\[
 P(N)\ge(1/2+\varepsilon)2^b
\]
has at most
\[
 K_{j,\varepsilon}(b+1)^{j-1}2^{b/2}
\]
nonempty edge sets that are unions of at most $j$ vertex-disjoint cycles.
For each fixed $j$, the power $j-1$ cannot be reduced throughout this class.
\end{theorem}

\begin{proof}
First suppose $N$ is bridgeless.
Use the distance layers and initial whole-layer marking from Lemma~\ref{lem:thin-layer-cycle-count}.
Additive stability bounds both $\Delta$ and $H$ by $C_\varepsilon$.
Thus at most $10C_\varepsilon+2$ vertices, and at most that many layers, are initially marked.
The maximum distance is $D\le b+2C_\varepsilon$.

Consider a union of exactly $m\le j$ vertex-disjoint cycles.
Write $I_1,\ldots,I_m$ for their intervals of distance layers.
Any layer strictly inside two of these intervals contains at least four vertices of their cycles.
Such a layer cannot be initially ordinary, since ordinary layers have only two vertices.
Only four layers can be endpoints of a given pair of intervals.
Consequently,
\[
 |I_a\cap I_b|\le B:=\lceil10C_\varepsilon+2\rceil+4
 \qquad(a\ne b).
\]

Fix the intervals, and mark all their endpoint layers.
This adds at most $4m$ vertices.
Fix all selected edges incident with marked vertices.
This costs a constant depending only on $j,\varepsilon$.
Every remaining ordinary layer is unused or carries the two through strands of one cycle.
The intervals determine this choice.
Across consecutive occupied ordinary layers, the selected edges form a perfect matching.
Two-choice cuts are nonconsecutive, so their number is at most $\tfrac12\sum_a|I_a|$.
These data determine the complete selected edge set.

Lemma~\ref{lem:fixed-interval-tuple-sum} bounds the number of unions with exactly $m$ components by
\[
 O_{j,\varepsilon}((D+1)^{m-1}2^{D/2})
 =O_{j,\varepsilon}((b+1)^{m-1}2^{b/2}).
\]
Sum over $1\le m\le j$ to prove the bridgeless case.

For a general normalized two-pole, its relevant bridgeless pieces occur in series.
Let their positive ranks be $b_1,\ldots,b_q$, with sum $b$.
Every piece retains normalized path density at least $1/2+\varepsilon$.
Suppose a packing uses $s$ pieces and has $m_i\ge1$ cycle components in piece $i$, with $\sum m_i\le j$.
The product of the polynomial factors is at most
\[
 \prod_{i\text{ used}}(b_i+1)^{m_i-1}
 \le(b+1)^{j-s}.
\]
For every fixed $s$, one has
\[
 \frac{\displaystyle\sum_{I\subseteq[q],\ |I|=s}
             2^{\frac12\sum_{i\in I}b_i}}
      {2^{b/2}}
 \le\binom qs2^{-(q-s)/2}\le c_s,
\]
where $c_s$ is finite because exponential decay dominates the fixed-degree polynomial in $q$.
There are only a bounded number of component-count allocations for fixed $j,s$.
Since $s\ge1$, the resulting polynomial degree is at most $j-1$.
Summing over $1\le s\le j$ proves the stated bound.

For sharpness, take the saturated capped ladder with $m$ square blocks.
Its rank is $b=2m+1$.
Partition its square blocks into $j$ consecutive groups, each containing at least two squares.
A group of $\ell$ squares has $2^\ell$ cycles that meet every square of the group.
There are two cap choices in each end square and two through choices in each interior square.
The cycles from different groups are vertex-disjoint.
Their block supports recover the chosen partition.
Thus this graph has at least
\[
 \binom{m-j-1}{j-1}2^m
 =\Omega_j(b^{j-1}2^{b/2})
\]
unions of exactly $j$ vertex-disjoint cycles, for $m\ge2j$.
This proves the optimality of the polynomial degree.
\end{proof}
